\section{Related Work and Evaluation Scope}
\paragraph{Serving mechanisms.}
Orca introduced iteration-level scheduling to change batches as requests
progress \citep{yu2022orca}. vLLM's PagedAttention addresses KV-memory
fragmentation and sharing \citep{kwon2023vllm}. Sarathi-Serve uses chunked
prefill to manage throughput--latency trade-offs \citep{agrawal2024sarathi}.
Bari et al. distinguish the throughput-optimal RAD construction from the
practical SLAI scheduler, which uses measured decode urgency and known
prompt lengths \citep{slai2025}. These are system and algorithm contributions
with their own execution assumptions. Our whole-FCFS, chunked-FCFS,
decode-first, and SLO-gated adaptations isolate a smaller set of admissible
rules. In particular, full-context reservations do not reproduce
PagedAttention, and an inactive urgency gate supplies no evidence for
SLAI's complete design or RAD's theorem. Disaggregated serving and dynamic
migration, exemplified by DistServe and Llumnix
\citep{zhong2024distserve,sun2024llumnix}, are outside the single-device,
nonpreemptive contract evaluated here.

\paragraph{Workload-sensitive evaluation.}
Papaioannou and Doudali explicitly examine how dataset and inference-scenario
choice affect comparisons among LLM inference systems
\citep{papaioannou2024workload}. Their work is a close precedent for the
importance of workload choice, rather than a gap this paper first identifies.
Etalon identifies limitations of conventional latency and throughput
metrics for streaming user experience and introduces a fluidity index
\citep{agrawal2024etalon}. Our metric comparisons use the frozen aggregate
endpoints; they neither introduce a new user-experience metric nor measure
streaming fluidity.
BurstGPT contributes production-derived workload characterization
\citep{burstgpt2025}, and ServeGen studies production workloads and their
generation \citep{servegen2026}. Workload resampling has a longer history
in parallel-job scheduler evaluation \citep{zakay2014resampling}.
Our analysis instead holds a selected trace corpus fixed, matches each
adaptation on the same window and condition, and distinguishes source
means from cross-window comparisons. We do not infer that the selected
windows are independent draws from provider populations.

\paragraph{Throughput optimality and finite-window completion.}
Dai, Deng, Li, and Peng's May 18, 2026 revision studies batching and
work-conserving scheduling through queueing stability
\citep{dai2026throughput}. Their token-budget model has batch time depending
on token load; their maximal-batch-size extension shows that request-count
constraints materially change the stability analysis. This makes the
binding constraint a substantive modeling choice, rather than an incidental
configuration detail. In \cm, full-context reservations limit residency to
four or eight requests, chunked work cannot exhaust the token budget, and
every batch costs the same time. Moreover, our throughput is output from a
finite, fully drained input divided by its completion duration. It is not
their long-run stability criterion. Their work motivates examining batching
and concurrency together, but its theorems do not certify these adaptations;
our finite-window throughput losses neither refute their stability results
nor establish a new scheduling principle.

\paragraph{Related manuscripts and incremental question.}
Related active manuscripts by the author already study inactive mechanisms
and metric sensitivity. The action-opportunity/headroom manuscript asks when
an executable alternative exists and measures reference-conditional,
one-step causal value. The module-counterfactual manuscript studies typed
transplants, downstream erasure and transplant fidelity. The closed-API
resource-accounting manuscript separates nominal budgets, realized resources
and answer selection on different tasks and data. M3 does not contribute
those methods or reuse their quantitative results as validation. Its
incremental object is a complete common-contract adaptation comparison:
exact trajectory consequences for the chunked/decode pair, measured remaining
trade-offs, and the contrast between source-mean stability and window/metric
rankings, with a separate historical-inference diagnostic. The shared ideas
of explicit accounting and conditional mechanisms are acknowledged, not
claimed as unique to this paper. Appendix~\ref{app:evidence} and the related-
manuscript disclosure identify inspected local versions; actual submitted
versions and public citation identities require author confirmation.
M3 also reuses windows, reconstruction and the archived 36/990 rule result
from the earlier portability line, whose rejected/withdrawn submissions are
not active concurrent work.

\paragraph{Serving simulators and model credibility.}
Vidur combines profiling and predictive models to explore LLM inference
configurations \citep{vidur2024}; LLMServingSim develops HW/SW co-simulation
for serving workloads \citep{llmservingsim2024}. LLMServingSim 2.0 further
models interactions among runtime scheduling, heterogeneous hardware, and
disaggregated resources \citep{llmservingsim2026}. Those frameworks target
hardware-aware performance estimation and evaluate prediction accuracy.
\cm\ uses fixed-duration batches and an explicit work ledger to make
adaptations comparable and execution independently checkable. It cannot
replace their predictive objectives or borrow their validation evidence.
Following the distinction between verification and validation in simulation
methodology \citep{sargent2013}, reconstructing an execution journal tests
whether the program implements the declared model. Establishing adequacy
for a physical serving system would require separate evidence.

\paragraph{Statistical procedures.}
The calibration diagnostic uses existing procedures: Benjamini--Hochberg
(BH), Benjamini--Yekutieli (BY), Holm, and bounded concentration inequalities
\citep{benjamini1995,benjamini2001dependency,holm1979,hoeffding1963}.
Their assumptions do not follow from successful code execution. In
particular, a multiplicity adjustment cannot establish validity of its
input p-values or repair a mismatch between the tested null and a practical
performance claim. STAT01 investigates those issues in specified synthetic
scenarios; it does not propose a replacement test.
